\section{Проектирование физического движка}

\subsection{Концептуальная модель физического мира}

Физический движок оперирует понятием \textit{физического мира} (\texttt{PhysicsWorld}) --- изолированного пространства симуляции, содержащего все объекты сцены. Мир является единственным владельцем всех создаваемых объектов: тел, коллайдеров и ограничений. Такая архитектура с единственной точкой владения (ownership) исключает висячие ссылки и упрощает сериализацию состояния для целей воспроизводимой симуляции.

Физический мир характеризуется следующими глобальными параметрами:
\begin{itemize}
  \item вектор гравитационного ускорения $\mathbf{g}$ (по умолчанию $(0, 0, -9{,}81)$\,м/с²);
  \item шаг интегрирования $\Delta t$ (фиксированный или переменный с ограничением сверху);
  \item число итераций решателя ограничений $n_{\text{iter}}$;
  \item флаг включения GPU-ускорения и число параллельных потоков.
\end{itemize}

\subsection{Представление твёрдых тел}
\label{sec:rigid_body}

Твёрдое тело (\texttt{RigidBody}) является центральным объектом физической симуляции. Его состояние полностью определяется набором атрибутов, перечисленных в таблице~\ref{tab:rigidbody_state}.

\begin{table}[h]
\centering
\caption{Атрибуты состояния объекта \texttt{RigidBody}}
\label{tab:rigidbody_state}
\begin{tabular}{lll}
\toprule
\textbf{Атрибут} & \textbf{Тип} & \textbf{Описание} \\
\midrule
\texttt{position}        & $\mathbb{R}^3$  & Положение центра масс в мировых координатах \\
\texttt{orientation}     & $\mathbb{H}^1$  & Единичный кватернион ориентации \\
\texttt{linear\_vel}     & $\mathbb{R}^3$  & Линейная скорость центра масс \\
\texttt{angular\_vel}    & $\mathbb{R}^3$  & Угловая скорость в мировых координатах \\
\texttt{inv\_mass}       & $\mathbb{R}$    & Обратная масса ($0$ для статических тел) \\
\texttt{inv\_inertia}    & $\mathbb{R}^{3\times3}$ & Обратный тензор инерции в локальных координатах \\
\texttt{force\_accum}    & $\mathbb{R}^3$  & Накопленная сила за текущий шаг \\
\texttt{torque\_accum}   & $\mathbb{R}^3$  & Накопленный момент за текущий шаг \\
\texttt{restitution}     & $[0,1]$         & Коэффициент упругости при ударе \\
\texttt{friction}        & $\mathbb{R}^+$  & Коэффициент трения Кулона \\
\texttt{is\_static}      & \texttt{bool}   & Признак статического (бесконечно массивного) тела \\
\texttt{is\_sleeping}    & \texttt{bool}   & Признак неактивного тела \\
\bottomrule
\end{tabular}
\end{table}

Обратная масса $m^{-1}$ используется вместо прямого хранения массы, так как это исключает деление в вычислениях горячего пути и позволяет естественно представить статические тела ($m^{-1} = 0$) без специальных проверок. Аналогично, для тензора инерции хранится обратная матрица в локальных координатах тела $\mathbf{I}_{\text{body}}^{-1}$; пересчёт в мировые координаты производится через матрицу вращения по формуле, рассмотренной в разделе~2.1.

\textit{Режим сна} (\textit{sleeping}) позволяет пропускать симуляцию тел, кинетическая энергия которых опустилась ниже заданного порога за несколько последовательных шагов. Тело переводится обратно в активное состояние при любом контакте с другим активным телом или при воздействии внешней силы. Данный механизм критически важен для производительности сцен с большим числом тел, большинство из которых в типичный момент времени находятся в покое.

\subsection{Коллайдеры и примитивы столкновений}
\label{sec:colliders}

Коллайдер (\texttt{Collider}) определяет геометрическую форму тела для целей обнаружения столкновений и, как правило, является упрощённым представлением визуальной геометрии. Движок поддерживает иерархию коллайдеров, показанную на рисунке~\ref{fig:collider_hierarchy}.

\begin{figure}[h]
\centering
\begin{tikzpicture}[
  font=\small,
  box/.style={draw, rounded corners=3pt, minimum width=3cm, minimum height=0.7cm,
              align=center, fill=white},
  basebox/.style={draw, rounded corners=3pt, minimum width=4cm, minimum height=0.7cm,
                  align=center, fill=gray!15},
  arr/.style={-{Stealth[length=5pt]}, thick},
]

\node[basebox] (base) at (0,0) {\texttt{Collider} (абстрактный)};

\node[box] (sphere)  at (-5.5,-1.7) {\texttt{SphereCollider}};
\node[box] (capsule) at (-2.5,-1.7) {\texttt{CapsuleCollider}};
\node[box] (box)     at (0.5,-1.7)  {\texttt{BoxCollider}};
\node[box] (cvx)     at (3.5,-1.7)  {\texttt{ConvexHullCollider}};
\node[box] (mesh)    at (6.5,-1.7)  {\texttt{TriMeshCollider}};

\draw[arr] (base.south) -- ++(-5.5,-0.6) -- (sphere.north);
\draw[arr] (base.south) -- ++(-2.5,-0.6) -- (capsule.north);
\draw[arr] (base.south) -- ++(0.5,-0.6)  -- (box.north);
\draw[arr] (base.south) -- ++(3.5,-0.6)  -- (cvx.north);
\draw[arr] (base.south) -- ++(6.5,-0.6)  -- (mesh.north);

\end{tikzpicture}
\caption{Иерархия типов коллайдеров физического движка}
\label{fig:collider_hierarchy}
\end{figure}

\textbf{SphereCollider} задаётся радиусом и смещением центра относительно центра масс тела. Это простейший примитив: алгоритм столкновения двух сфер сводится к сравнению расстояния между центрами с суммой радиусов.

\textbf{CapsuleCollider} --- цилиндр с полусферическими торцами, задаётся радиусом и длиной оси. Хорошо аппроксимирует звенья манипуляторов, конечности роботов и трубообразные конструкции при низкой стоимости проверки столкновений.

\textbf{BoxCollider} (ориентированный параллелепипед, OBB) задаётся тремя полуразмерами вдоль осей локальной системы координат тела. Эффективен для прямоугольных конструкций.

\textbf{ConvexHullCollider} представляет выпуклую оболочку набора точек. Алгоритм GJK, рассмотренный в разделе~2.2, работает с любым выпуклым телом через функцию опоры; выпуклая оболочка является наиболее точным приближением произвольной геометрии, для которого GJK гарантированно применим.

\textbf{TriMeshCollider} --- произвольная невыпуклая сетка треугольников. Используется только для статических тел (пол, стены, неподвижные препятствия), так как столкновение двух произвольных невыпуклых тел требует разбиения на выпуклые части. Проверка столкновения с динамическим телом производится перебором по треугольникам с ускорением через BVH (\textit{Bounding Volume Hierarchy}).

Каждый коллайдер хранит смещение (\texttt{offset}) и ориентацию (\texttt{offset\_rotation}) относительно центра масс тела, что позволяет прикреплять несколько коллайдеров к одному телу и точнее аппроксимировать сложную геометрию.

\subsection{Ограничения и сочленения}
\label{sec:constraints}

Ограничения (\textit{constraints}) реализуют связи между парами тел. В рамках данного раздела рассмотрена классификация и структура данных; математическая формулировка дана в разделе~2.4.

Базовый класс \texttt{Constraint} задаёт интерфейс, используемый решателем:

\begin{lstlisting}[language=C++, caption={Интерфейс базового класса ограничения}]
class Constraint {
public:
    RigidBody *bodyA, *bodyB;  // nullptr для мирового фрейма

    // Вычислить якобиан J и вектор смещения b для данного шага
    virtual void computeJacobian(JacobianBlock &J, float b[],
                                 float dt) const = 0;

    // Нижняя и верхняя границы импульса (для трения и упоров)
    virtual void impulseLimits(float lo[], float hi[],
                               const float lambda[]) const = 0;
};
\end{lstlisting}

Конкретные классы ограничений, реализуемые в движке:

\begin{itemize}
  \item \texttt{ContactConstraint} --- создаётся автоматически при обнаружении столкновения: одно нормальное и до двух тангенциальных (фрикционных) ограничений;
  \item \texttt{RevoluteJoint} --- вращательная пара: фиксирует 5 из 6 степеней свободы, оставляя одну ось вращения; поддерживает пределы угла и мотор с заданной скоростью/моментом;
  \item \texttt{PrismaticJoint} --- поступательная пара вдоль одной оси;
  \item \texttt{FixedJoint} --- жёсткое соединение (6 ограничений), применяется для объединения нескольких тел в составное;
  \item \texttt{BallJoint} --- шаровое сочленение (фиксирует 3 трансляционные степени свободы);
  \item \texttt{DistanceConstraint} --- поддерживает фиксированное или ограниченное расстояние между двумя точками тел.
\end{itemize}

\subsection{Жизненный цикл шага симуляции}
\label{sec:sim_step}

Каждый шаг симуляции длительностью $\Delta t$ выполняется в строго определённой последовательности, схематически представленной на рисунке~\ref{fig:sim_pipeline}.

\begin{figure}[h]
\centering
\begin{tikzpicture}[
  font=\small,
  simstep/.style={draw, rounded corners=3pt, fill=white,
               minimum width=8cm, minimum height=0.75cm, align=center},
  gpustep/.style={draw, rounded corners=3pt, fill=blue!10,
              minimum width=8cm, minimum height=0.75cm, align=center},
  arr/.style={-{Stealth[length=6pt]}, thick},
  note/.style={font=\footnotesize\itshape, text=gray, align=left},
]

\node[simstep] (forces)   at (0,0)    {1. Применение внешних сил и гравитации};
\node[gpustep] (broad)    at (0,-1.1) {2. Широкая фаза обнаружения столкновений (GPU)};
\node[gpustep] (narrow)   at (0,-2.2) {3. Узкая фаза: GJK/EPA (GPU, параллельно по парам)};
\node[simstep] (warmstart)at (0,-3.3) {4. Warm-starting: перенос импульсов с прошлого шага};
\node[gpustep] (solver)   at (0,-4.4) {5. Итерационный решатель PGS (GPU)};
\node[simstep] (integrate)at (0,-5.5) {6. Симплектическое интегрирование положений};
\node[simstep] (sleep)    at (0,-6.6) {7. Обновление режима сна тел};
\node[simstep] (sync)     at (0,-7.7) {8. Передача трансформаций в слой сцены};

\foreach \a/\b in {forces/broad, broad/narrow, narrow/warmstart,
                    warmstart/solver, solver/integrate,
                    integrate/sleep, sleep/sync}
  \draw[arr] (\a.south) -- (\b.north);

\node[note, right=0.5cm of broad]    {§\,3.3};
\node[note, right=0.5cm of narrow]   {§\,3.3};
\node[note, right=0.5cm of solver]   {§\,3.3};

\end{tikzpicture}
\caption{Последовательность вычислений в рамках одного шага симуляции (голубым выделены стадии с GPU-ускорением)}
\label{fig:sim_pipeline}
\end{figure}

\textbf{Стадия 1.} Для каждого активного тела интегратор вычисляет предварительную скорость $\mathbf{v}^* = \mathbf{v} + m^{-1}\mathbf{F}_{\text{ext}}\Delta t$, где $\mathbf{F}_{\text{ext}}$ включает гравитацию и любые силы, приложенные пользовательским кодом. Эта предварительная скорость используется в качестве начального приближения для решателя.

\textbf{Стадии 2--3.} Обнаружение столкновений выполняется на GPU в два этапа (подробно рассмотрено в разделах~2.2 и~3.3). На выходе формируется список \textit{контактных точек} --- структур, содержащих нормаль, глубину проникновения, опорные точки на каждом теле и идентификаторы пары тел.

\textbf{Стадия 4.} Warm-starting (тёплый запуск) использует импульсы $\lambda$ из предыдущего шага в качестве начального приближения для итерационного решателя. Для каждого контакта, идентифицированного по хэшу пары (bodyA, bodyB, featureId), накопленный импульс немедленно применяется к скоростям тел, сокращая необходимое число итераций решателя.

\textbf{Стадия 5.} Решатель PGS (рассмотренный в разделе~2.4) выполняет $n_{\text{iter}}$ итераций; на GPU каждая итерация обрабатывает все независимые ограничения параллельно (подробнее в разделе~3.3).

\textbf{Стадия 6.} Симплектическое интегрирование Эйлера обновляет положения и ориентации:
\begin{align}
  \mathbf{r}(t + \Delta t) &= \mathbf{r}(t) + \mathbf{v}(t + \Delta t)\,\Delta t, \\
  \mathbf{q}(t + \Delta t) &= \text{norm}\!\left(\mathbf{q}(t) + \tfrac{1}{2}\boldsymbol{\Omega}(\boldsymbol{\omega}(t+\Delta t))\,\mathbf{q}(t)\,\Delta t\right).
\end{align}

\textbf{Стадия 7.} Энергетическое пороговое значение $E_{\text{sleep}} = \frac{1}{2}(m|\mathbf{v}|^2 + \boldsymbol{\omega}^T\mathbf{I}\boldsymbol{\omega})$ сравнивается с порогом; при превышении счётчик неактивности сбрасывается.

\textbf{Стадия 8.} Матрицы преобразования всех тел записываются в VSG-узлы сцены. Эта операция выполняется на CPU и является единственной точкой синхронизации между физическим движком и подсистемой рендеринга.

\subsection{Хранение данных в памяти}

Для эффективного GPU-вычисления тела и их атрибуты хранятся в структурах, ориентированных на массивный параллелизм. Вместо объектно-ориентированного представления «одна структура --- одно тело» используется шаблон \textit{Structure of Arrays} (SoA): каждый атрибут всех тел хранится в отдельном непрерывном массиве.

\begin{lstlisting}[language=C++, caption={Структура хранения данных тел (SoA-формат)}]
struct RigidBodyData {
    // Кинематические данные
    std::vector<glm::vec3>  positions;       // [N]
    std::vector<glm::quat>  orientations;    // [N]
    std::vector<glm::vec3>  linear_vels;     // [N]
    std::vector<glm::vec3>  angular_vels;    // [N]

    // Инерционные данные (загружаются один раз)
    std::vector<float>      inv_masses;      // [N]
    std::vector<glm::mat3>  inv_inertias;    // [N]

    // Приложенные силы (сбрасываются каждый шаг)
    std::vector<glm::vec3>  force_accums;    // [N]
    std::vector<glm::vec3>  torque_accums;   // [N]

    // Флаги состояния (упакованы побитово)
    std::vector<uint32_t>   flags;           // [N/32]
};
\end{lstlisting}

SoA-формат обеспечивает последовательный доступ к памяти при обработке массива тел, что критично как для CPU-кэша, так и для GPU-вейвфронтов. При передаче данных в GPU все массивы копируются в единый~SSBO (Shader Storage Buffer Object), разметка которого подробно рассмотрена в разделе~3.3.
